\chapter*{怎样读这一卷}
\addcontentsline{toc}{chapter}{怎样读这一卷}

本卷从一六四四年清军入关、北京政权建立写到一九一二年宣统帝退位。南明、三藩、郑氏台湾、准噶尔、革命与地方独立只写入与帝国形成或终结直接相连的部分；满洲、蒙古、新疆、青藏高原、东北、西南、台湾、海疆和海外交往都有自己的行动者与材料节奏，不能缩成单一“中原王朝”的年表。

先依清帝年号、干支、月日和换算建立年序；遇到南明、郑氏、朝鲜、日本、俄历、格里历、伊斯兰历、藏历或地方纪年时，应把并列日期看作核验而非多余细节。再区分长期的财政、交通、军事和区域条件，临近的触发，可用的文书、补给与联盟，以及结果怎样改变下一轮选择；条约、改革或战报本身不自动完成这种因果链。

《清实录》《清史稿》、奏折、军机处和内阁档、地方档案、报刊、条约、外文档案与满蒙藏维文材料提供不同视角。上谕、章程和条约能证明文书与规范，不能保证地方执行；户口、田赋、海关、军额、战果和灾荒数字须连同机构、地域、单位与口径阅读。官方或外文的群体标签也不是群体本质。

\texttt{event} 让读者进入首次深描的节点，\texttt{turningpoint} 标示路径转折，\texttt{causalchain} 拆开条件、触发与反馈；\texttt{textwindow} 限定原文所能说明的范围；\texttt{sourcenote}、\texttt{debate} 保留证据的距离、立场和分歧。请用这些停靠点回到多区域的连续叙事，而不要把它们读成已经替所有行动者作出的裁决。
